\subsection{A universal category determined by the three short products}

\begin{theorem}[Proved: a universal short-product obstruction]
\label{A7:universal-short-model}
Let $\mathcal C$ be a finite acyclic category with five objects and
full support $0<1<2<3<4$. Write its adjacent hom-sets as $A,B,C,D$
and its next three hom-sets as
\[
 U=\mathcal C(0,2),\qquad V=\mathcal C(1,3),
 \qquad W=\mathcal C(2,4).
\]
Let the chronological short-product maps be
$P:A\times B\to U$, $Q:B\times C\to V$ and $R:C\times D\to W$.
There is a finite acyclic category $F(P,Q,R)$, defined below solely
from these seven sets and three maps, such that
\[
 \dim_{\mathbb F_2}H_3(B\mathcal C;\mathbb F_2)
 \ \geq\ \dim_{\mathbb F_2}H_3(BF;\mathbb F_2).
\]
If every degree-two square of $BF$ over $\mathbb F_2$ is zero,
the same is true for $B\mathcal C$. Consequently either
$H_3(BF;\mathbb F_2)\ne0$ or $H^2(BF;\mathbb F_2)=0$
excludes $B\mathcal C\simeq\mathbb{CP}^2$.
\end{theorem}
\begin{proof}
Form two finite bipartite graphs. The first has tagged vertex set
$(A\times V)\sqcup(U\times C)$ and one edge for each $(a,b,c)$,
joining
\[
 (a,Q(b,c))\quad\text{to}\quad(P(a,b),c).
\]
Let $L$ be its set of connected components. The second has tagged
vertex set $(B\times W)\sqcup(V\times D)$, with an edge for each
$(b,c,d)$ joining
\[
 (b,R(c,d))\quad\text{to}\quad(Q(b,c),d).
\]
Let $K$ be its component set. Define $F$ to have the same five
objects, adjacent hom-sets $A,B,C,D$, short hom-sets $U,V,W$,
hom-sets $F(0,3)=L$, $F(1,4)=K$, and a singleton $F(0,4)$.
Its three short compositions are $P,Q,R$. A product in $A\times V$
or $U\times C$ is its graph component in $L$; the two analogous
products into $K$ are defined similarly. All products into $F(0,4)$
are the unique arrow there. Add the five identity arrows.

These rules define all products. Associativity on the object
strings $0123$ and $1234$ is exactly the relation imposed by the
two graphs. Every other triple of cross arrows starts at $0$ and
ends at $4$, so both associative composites are the unique arrow.
Thus $F$ is an actual finite acyclic category. The construction
includes every vertex of the two graphs, even if some element of
$U,V,W$ is not in the image of its short-product map.

First collapse $\mathcal C(0,4)$ to a singleton, obtaining
$\mathcal C_0$. By Proposition~\ref{A7:extreme-arrow-collapse},
the induced map is surjective on second cohomology and an
isomorphism on fourth cohomology. Its normalized chain groups in
degrees at least two, and all boundaries in degrees at least
three, are unchanged. In particular its third homology is
unchanged. We do not assert that this collapse preserves the
projective-plane homotopy type or the rank of second cohomology.

Call an arrow $0\to3$ generated if it is a product from
$A\times V$ or $U\times C$, and define generated arrows $1\to4$
using $B\times W$ or $V\times D$. Fold every nongenerated
$0\to3$ arrow onto a fixed generated one. This is an idempotent
endomorphism of the connecting profunctor across the cut
$\{0\}\mid\{1,2,3,4\}$. No two-arrow product hits a folded
arrow, and every further product lands in the singleton
$\mathcal C_0(0,4)$, so every naturality equation holds.
Likewise fold the nongenerated $1\to4$ arrows across the final
object cut. These operations give actual retractions onto the
subcategory $Y$ containing exactly the generated long arrows.

For either fold, write $i:Y'\hookrightarrow X$ and $r:X\to Y'$
for its inclusion and retraction. The relative mapping-cylinder
quotient is a suspension, as in the proof of
Theorem~\ref{CP2:A-cut-retract}, so all positive-degree relative
cup products vanish. For $\alpha\in H^2(BX;\mathbb F_2)$ put
$\lambda=\alpha-r^*i^*\alpha$. Since $i^*\lambda=0$,
$\lambda$ lifts to relative cohomology and $\lambda^2=0$.
Graded commutativity in characteristic two now gives
\[
 \alpha^2=r^*((i^*\alpha)^2).
\]
Thus zero degree-two squares transfer back across each fold
without a hypothesis on the dimension of second cohomology.

No arrow $0\to3$ or $1\to4$ can be a factor of a normalized
simplex of degree at least three. Moreover, a two-simplex
containing a nongenerated long arrow never occurs in a third
boundary: an inner face producing that long arrow would make
it generated, and an outer face would require such an arrow
already to be a factor of the three-simplex. Deleting these
arrows therefore deletes only zero rows from the third boundary;
the third and fourth chain groups and fourth boundary are
identical. Hence $H_3(BY;\mathbb F_2)=H_3(B\mathcal C_0;\mathbb F_2)$.

Associativity in $Y$ gives an actual functor $F\to Y$ which is
the identity on every adjacent and short hom-set. Its maps on
$L,K$ send graph components to their common long composite,
and are surjective because every long arrow of $Y$ is generated.
It is the identity on normalized chain groups in degrees three
and four, and gives a surjection $q_2:C_2(BF;\mathbb F_2)
\to C_2(BY;\mathbb F_2)$. The identity
\[
 \partial_3^Y=q_2\partial_3^F
\]
implies $\operatorname{rank}\partial_3^Y\leq
\operatorname{rank}\partial_3^F$. Since $\partial_4$ is identical,
this proves the asserted third-homology inequality.
The same identity on $C_3,C_4,\partial_4$ gives an isomorphism
on fourth cohomology. Naturality of the cup product then shows
that zero squares on $BF$ imply zero squares on $BY$. Transfer
back through the two folds, and then through the surjection
$H^2(B\mathcal C_0;\mathbb F_2)\to
H^2(B\mathcal C;\mathbb F_2)$, proves the square assertion.
The projective plane has zero third homology and a nonzero
degree-two square over $\mathbb F_2$, proving the obstruction.
\end{proof}

\begin{remark*}[Computational use]
The two component graphs can be computed before assigning any
long-product table. The theorem therefore allows a single check
of $F(P,Q,R)$ to exclude every longer-table completion of the
same three short maps. This is a proved reduction; it does not
by itself claim that any new finite search has been completed.
\end{remark*}
